\documentclass[10pt,a4paper]{amsart}
\usepackage{amsmath,amssymb,amsthm,graphicx}
\usepackage{verbatim}
\newenvironment{reference}{\comment}{\endcomment}
%\newenvironment{reference}{}{}
\newenvironment{slogan}{\comment}{\endcomment}
\newenvironment{history}{\comment}{\endcomment}

% For commutative diagrams we use Xy-pic
\usepackage[all]{xy}

% We use 2cell for 2-commutative diagrams.
\xyoption{2cell}
% Theorem environments.
%
\theoremstyle{plain}
\newtheorem{theorem}[subsection]{Theorem}
\newtheorem{proposition}[subsection]{Proposition}
\newtheorem{lemma}[subsection]{Lemma}

\theoremstyle{definition}
\newtheorem{definition}[subsection]{Definition}
\newtheorem{example}[subsection]{Example}
\newtheorem{exercise}[subsection]{Exercise}
\newtheorem{situation}[subsection]{Situation}

\theoremstyle{remark}
\newtheorem{remark}[subsection]{Remark}
\newtheorem{remarks}[subsection]{Remarks}

\numberwithin{equation}{subsection}

% Macros
%
\def\lim{\mathop{\mathrm{lim}}\nolimits}
\def\colim{\mathop{\mathrm{colim}}\nolimits}
\def\Spec{\mathop{\mathrm{Spec}}}
\def\Hom{\mathop{\mathrm{Hom}}\nolimits}
\def\Ext{\mathop{\mathrm{Ext}}\nolimits}
\def\SheafHom{\mathop{\mathcal{H}\!\mathit{om}}\nolimits}
\def\SheafExt{\mathop{\mathcal{E}\!\mathit{xt}}\nolimits}
\def\Sch{\mathit{Sch}}
\def\Mor{\mathop{\mathrm{Mor}}\nolimits}
\def\Ob{\mathop{\mathrm{Ob}}\nolimits}
\def\Sh{\mathop{\mathit{Sh}}\nolimits}
\def\NL{\mathop{N\!L}\nolimits}
\def\CH{\mathop{\mathrm{CH}}\nolimits}
\def\proetale{{pro\text{-}\acute{e}tale}}
\def\etale{{\acute{e}tale}}
\def\QCoh{\mathit{QCoh}}
\def\Ker{\mathop{\mathrm{Ker}}}
\def\Im{\mathop{\mathrm{Im}}}
\def\Coker{\mathop{\mathrm{Coker}}}
\def\Coim{\mathop{\mathrm{Coim}}}

% Boxtimes
%
\DeclareMathSymbol{\boxtimes}{\mathbin}{AMSa}{"02}

%
% Macros for moduli stacks/spaces
%
\def\QCohstack{\mathcal{QC}\!\mathit{oh}}
\def\Cohstack{\mathcal{C}\!\mathit{oh}}
\def\Spacesstack{\mathcal{S}\!\mathit{paces}}
\def\Quotfunctor{\mathrm{Quot}}
\def\Hilbfunctor{\mathrm{Hilb}}
\def\Curvesstack{\mathcal{C}\!\mathit{urves}}
\def\Polarizedstack{\mathcal{P}\!\mathit{olarized}}
\def\Complexesstack{\mathcal{C}\!\mathit{omplexes}}
% \Pic is the operator that assigns to X its picard group, usage \Pic(X)
% \Picardstack_{X/B} denotes the Picard stack of X over B
% \Picardfunctor_{X/B} denotes the Picard functor of X over B
\def\Pic{\mathop{\mathrm{Pic}}\nolimits}
\def\Picardstack{\mathcal{P}\!\mathit{ic}}
\def\Picardfunctor{\mathrm{Pic}}
\def\Deformationcategory{\mathcal{D}\!\mathit{ef}}
\usepackage{cleveref}
\crefname{originaldefinition}{original definition}{original definitions}
\Crefname{originaldefinition}{Original definition}{Original definitions}
\makeatletter
\newlabel{definition-relative-global-complete-intersection}{{00SP}{source}{}{sourcepointer-1}{}}
\newlabel{definition-relative-global-complete-intersection@cref}{{[originaldefinition][0][]00SP}{[1][0][]source}{}{}{}}
\makeatother
\crefname{originallemma}{original lemma}{original lemmas}
\Crefname{originallemma}{Original lemma}{Original lemmas}
\makeatletter
\newlabel{lemma-dimension-fibres-bounded-open-upstairs}{{00QH}{source}{}{sourcepointer-2}{}}
\newlabel{lemma-dimension-fibres-bounded-open-upstairs@cref}{{[originallemma][0][]00QH}{[1][0][]source}{}{}{}}
\makeatother
\crefname{originallemma}{original lemma}{original lemmas}
\Crefname{originallemma}{Original lemma}{Original lemmas}
\makeatletter
\newlabel{lemma-dimension-preserved-field-extension}{{00P3}{source}{}{sourcepointer-3}{}}
\newlabel{lemma-dimension-preserved-field-extension@cref}{{[originallemma][0][]00P3}{[1][0][]source}{}{}{}}
\makeatother
\title{A relative global complete intersection is defined over a finitely generated subring of the base}
\author{The Stacks Project Authors}
\date{stacks-project:algebra.tex:local-snapshot-20261012; GNU FDL 1.2 or later}
\begin{document}
\maketitle
\noindent\textit{Editorial note.} All delivered mathematical text is the author's own, copied from the pinned Stacks Project chapter file algebra.tex; no statement or proof step was rewritten, reordered or removed.\par\medskip
\noindent\textit{Editorial note.} Preamble provenance, as required for this source: the frozen snapshot stages the chapter's own preamble file beside it (preamble.tex, 217 lines, 8116 bytes, the same shared Stacks-wide preamble the published chow and etale-cohomology cards use); the first line of the pinned chapter file is \textbackslash{}input\{preamble\}, so no preamble text lives inside the chapter itself. The card takes that file's lines 8-18 (the verbatim package and the three \textbackslash{}comment-based environments reference/slogan/history at lines 9-12, section 13.2, together with the Xy-pic setup) and its lines 154 to the runtime end of the file (the native theorem declarations and mathematical macros). Not one character was changed, no package was added, no macro was defined, and the excerpt serves only as preamble, never as a span coordinate.\par\medskip
\noindent\textit{Editorial note.} The same-chapter pointers are located at their \textbackslash{}label lines inside the pinned chapter file; the cross-chapter pointers are located in the project's own tag list staged in the snapshot (sources/stacks-chow/tags.txt), at the verbatim lines carrying their canonical tags.\par\medskip
\noindent\textit{Editorial note.} The counted statement refers to the definition of a relative global complete intersection; it is recorded as a navigation stub with kind definition, and the arguments using it stay inside the counted span.\par\medskip
\noindent\textit{Editorial note.} The counted proof also refers to same-chapter statements by label, recorded the same way; no \textbackslash{}cite appears, so the card carries no bibliography.\par\medskip
\noindent\textit{Allowed standard prerequisites (editorial).}\begin{itemize}
\item Finitely many coefficients generate a subring. Used at the start of the counted proof (lines 36957-36959). Produces the finitely generated base over which the polynomials are already defined.
\item Finiteness of the number of primes where a condition fails. Used in the counted proof at lines 36960-36980. Shows the defining conditions hold over the small base as well.
\end{itemize}\medskip
\ifdefined\hypertarget\hypertarget{sourcepointer-1}{}\fi
\noindent\textit{Source reference.} Original definition 00SP: sources/stacks-algebra/pinned-original-source.tex, source lines 36723--36723. Navigation pointer for \textbackslash{}ref\{definition-relative-global-complete-intersection\} used inside the counted span (line 36947); the argument that consumes it is retained in the span.\par\medskip
\ifdefined\hypertarget\hypertarget{sourcepointer-2}{}\fi
\noindent\textit{Source reference.} Original lemma 00QH: sources/stacks-algebra/pinned-original-source.tex, source lines 31525--31525. Navigation pointer for \textbackslash{}ref\{lemma-dimension-fibres-bounded-open-upstairs\} used inside the counted span (line 36967); the argument that consumes it is retained in the span.\par\medskip
\ifdefined\hypertarget\hypertarget{sourcepointer-3}{}\fi
\noindent\textit{Source reference.} Original lemma 00P3: sources/stacks-algebra/pinned-original-source.tex, source lines 28507--28507. Navigation pointer for \textbackslash{}ref\{lemma-dimension-preserved-field-extension\} used inside the counted span (line 36966); the argument that consumes it is retained in the span.\par\medskip
\begin{lemma}
\label{lemma-relative-global-complete-intersection-Noetherian}
Let $R$ be a ring. Let $S = R[x_1, \ldots, x_n]/(f_1, \ldots, f_c)$
be a relative global complete intersection
(Definition \ref{definition-relative-global-complete-intersection}).
There exist a finite type $\mathbf{Z}$-subalgebra $R_0 \subset R$
such that $f_i \in R_0[x_1, \ldots, x_n]$ and such that
$$
S_0 = R_0[x_1, \ldots, x_n]/(f_1, \ldots, f_c)
$$
is a relative global complete intersection.
\end{lemma}

\begin{proof}
Let $R_0 \subset R$ be the $\mathbf{Z}$-algebra of $R$ generated by all the
coefficients of the polynomials $f_1, \ldots, f_c$. Let
$S_0 = R_0[x_1, \ldots, x_n]/(f_1, \ldots, f_c)$.
Clearly, $S = R \otimes_{R_0} S_0$.
Pick a prime $\mathfrak q \subset S$ and denote
$\mathfrak p \subset R$, $\mathfrak q_0 \subset S_0$, and
$\mathfrak p_0 \subset R_0$ the primes it lies over.
Because $\dim (S \otimes_R \kappa(\mathfrak p) ) = n - c$
we also have $\dim (S_0 \otimes_{R_0} \kappa(\mathfrak p_0)) = n - c$,
see Lemma \ref{lemma-dimension-preserved-field-extension}.
By Lemma \ref{lemma-dimension-fibres-bounded-open-upstairs}
there exists a $g \in S_0$, $g \not \in \mathfrak q_0$
such that all nonempty fibres of $R_0 \to (S_0)_g$
have dimension $\leq n - c$. As $\mathfrak q$ was arbitrary and
$\Spec(S)$ quasi-compact, we can find finitely many
$g_1, \ldots, g_m \in S_0$ such that (a) for $j = 1, \ldots, m$
the nonempty fibres of
$R_0 \to (S_0)_{g_j}$ have dimension $\leq n - c$ and (b) the image of
$\Spec(S) \to \Spec(S_0)$ is contained in $D(g_1) \cup \ldots \cup D(g_m)$.
In other words, the images of $g_1, \ldots, g_m$ in $S = R \otimes_{R_0} S_0$
generate the unit ideal. After increasing $R_0$ we may assume
that $g_1, \ldots, g_m$ generate the unit ideal in $S_0$. By (a)
the nonempty fibres of $R_0 \to S_0$ all have dimension $\leq n - c$
and we conclude.
\end{proof}

\end{document}
